\chapter{Fundamentação Teórica}\label{cap:fundTeoric}
\addcontentsline{toc}{chapter}{Fundamentação Teórica}
Este capítulo estabelece a base conceitual sobre a qual o experimento deste trabalho é construído. 

A Seção ~\ref{sec:seguranca-aplicacoes} trata da segurança de aplicações e da análise estática: parte dos conceitos gerais de vulnerabilidade e das taxonomias CWE/CVE, avança para o funcionamento, as vantagens e as limitações estruturais do \textit{Static Application Security Testing} (SAST) e conclui com a descrição das duas ferramentas SAST  (Bandit e Semgrep Community Edition). 

A Seção ~\ref{sec:llms} apresenta os Modelos de Linguagem de Grande Escala (LLMs): seus fundamentos arquiteturais, sua aplicação ao código-fonte e o mecanismo de \textit{Retrieval-Augmented Generation} (RAG), relevante para entender como ferramentas de codificação recuperam contexto de um repositório. 

A Seção ~\ref{sec:agentes-codificacao} discute os agentes de codificação assistidos por IA - a evolução do paradigma de autocomplete para agentes autônomos, o que caracteriza tecnicamente um agente, as ferramentas empresariais avaliadas (Cursor e OpenAI Codex) e o uso específico de IA para segurança de código, incluindo abordagens híbridas com SAST. 

Por fim, a Seção ~\ref{sec:metricas-avaliacao} define as métricas quantitativas e qualitativas empregadas para avaliar ferramentas de detecção de vulnerabilidades, encerrando com uma contextualização dos benchmarks utilizados na literatura recente, detalhados no Capítulo ~\ref{cap:estadoarte}.


\section{Segurança de Aplicações e Análise Estática}\label{sec:seguranca-aplicacoes}
\subsection{Vulnerabilidades de software: conceitos gerais}\label{subsec:vulnerabilidades-software}
Uma vulnerabilidade de software é uma falha na especificação, no projeto ou na implementação de um sistema que pode ser explorada para comprometer sua confidencialidade, integridade ou disponibilidade \cite{MITRECWE2024}. Essa definição situa o problema em dois planos distintos, msas complementares: o plano da \textit{fraqueza} (weakness) , a característica estrutural do código que torna a exploração possível, por exemplo a ausência de sanitização de uma entrada externa antes de seu uso em uma consulta SQL  e o plano da \textit{vulnerabilidade concreta} (vulnerability), a instância real dessa fraqueza em um produto de software específico, potencialmente já documentada e associada a um identificador público. É justamente essa distinção que estrutura as duas taxonomias descritas na Seção  ~\ref{subsec:cwe-cve}.

A motivação para detectar vulnerabilidades o mais cedo possível no ciclo de desenvolvimento é praticamente consensual na literatura, tanto pelo custo crescente de correção em estágios avançados quanto pelo risco de exploração em produção \cite{Dalaq2025SASTSLR, Shiri2024MLSLR}. Esse princípio conhecido como \textit{shift-left}  motivou o movimento \textbf{DevSecOps}, que propõe integrar atividades de segurança (incluindo análise estática automatizada) diretamente aos pipelines de integração e entrega contínuas (CI/CD), em vez de tratá-las como uma fase de auditoria isolada ao final do desenvolvimento \cite{Myrbakken2017DevSecOps}. A adoção de ferramentas de análise automatizada nas fases iniciais do ciclo de vida do software é, portanto, uma consequência direta dessa mudança de postura \cite{Wadhams2024SASTBarriers}.

O problema ganhou uma nova dimensão com a popularização de assistentes de programação baseados em IA generativa. Um estudo pioneiro sobre o GitHub Copilot verificou que aproximadamente 40\% dos programas gerados em cenários alinhados ao \textit{MITRE Top 25 CWE} continham vulnerabilidades \cite{Pearce2022CopilotSecurity}. De forma complementar, um experimento controlado com usuários demonstrou que desenvolvedores assistidos por IA produziram código menos seguro do que um grupo de controle e, paradoxalmente, relataram maior confiança na segurança do próprio resultado \cite{Perry2023InsecureAICode}. Esses achados indicam uma tendência de crescimento do volume de código potencialmente vulnerável em produção, o que reforça a relevância de mecanismos automatizados tradicionais ou baseados em IA capazes de detectar essas falhas antes da implantação, problema central deste trabalho.

\subsection{CWE (Common Weakness Enumeration) e CVE (Common Vulnerabilities and Exposures)}\label{subsec:cwe-cve}
Para viabilizar a comunicação, a comparação de ferramentas e a construção de conjuntos de dados rotulados, a comunidade de segurança de software adota duas taxonomias mantidas pela MITRE Corporation, complementares entre si. A \textbf{Common Weakness Enumeration (CWE)} é um dicionário de classes abstratas de fraquezas de software, organizado hierarquicamente: cada entrada CWE descreve um \textit{tipo} de falha , por exemplo, CWE-89 (\textit{SQL Injection}), CWE-22 (\textit{Path Traversal}) ou CWE-78 (\textit{OS Command Injection}) independentemente de qualquer produto específico em que ela ocorra \cite{MITRECWE2024}. Já o registro \textbf{Common Vulnerabilities and Exposures (CVE)} cataloga ocorrências concretas de vulnerabilidades, cada uma identificada por um código único (por exemplo, CVE-2024-XXXXX) e tipicamente associada a um ou mais identificadores CWE que descrevem sua natureza \cite{MITRECVE2024}. Em outras palavras, a CWE responde "que tipo de falha é essa?", enquanto a CVE responde "onde e quando essa falha específica ocorreu?".

Essas taxonomias cumprem três papéis neste trabalho. Primeiro, servem de vocabulário comum para categorizar e comparar os achados das três abordagens avaliadas SAST, ferramentas empresariais de IA e a combinação híbrida , permitindo verificar se a cobertura por categoria de vulnerabilidade é equivalente entre elas, um cuidado metodológico recomendado por estudos que apontam que nenhuma ferramenta domina todas as categorias CWE simultaneamente \cite{Esposito2024SASTComparison, Zhu2024SASTAndroid}. Segundo, orientam a curadoria de amostras rotuladas de código Python vulnerável utilizadas no experimento, uma vez que datasets como o VUDENC são construídos precisamente a partir de correções de commits associadas a tipos CWE específicos. Terceiro, complementam-se com classificações de risco de mais alto nível voltadas a aplicações web, como o \textbf{OWASP Top 10}, adotado como base de diversos benchmarks sintéticos citados na literatura de ferramentas híbridas (Seção ~\ref{subsubsec:hibrido-sast-ia}), entre eles o OWASP Benchmark \cite{OWASP2021Top10}. É importante notar, contudo, que essas taxonomias apresentam desbalanceamento de cobertura entre linguagens: mapeamentos sistemáticos da literatura mostram forte concentração de estudos e datasets rotulados em C/C++ e Java, com Python historicamente sub-representado \cite{Kalouptsoglou2023VulnPrediction, Germano2025LLMSLR} lacuna discutida com mais profundidade no Capítulo ~\ref{cap:estadoarte}.


\subsection{Static Application Security Testing (SAST)}\label{subsec:sast}
\subsubsection{Funcionamento e técnicas}\label{subsubsec:sast-funcionamento}
\textit{Static Application Security Testing} (SAST) designa a família de técnicas e ferramentas que analisam o código-fonte, o bytecode ou representações intermediárias de um programa em busca de padrões associados a vulnerabilidades, sem executar o programa \cite{Dalaq2025SASTSLR}. Essa ausência de execução é a característica definidora do paradigma e o distingue de técnicas dinâmicas (DAST, testes de penetração) e de abordagens híbridas de análise (IAST).

Do ponto de vista técnico, as ferramentas SAST tipicamente constroem uma ou mais representações intermediárias do programa  árvores sintáticas abstratas (\textit{Abstract Syntax Trees}, AST), grafos de fluxo de controle (\textit{Control Flow Graphs}) e grafos de fluxo de dados (\textit{Data Flow Graphs})  e aplicam sobre elas regras ou consultas (\textit{queries}) que codificam padrões considerados inseguros \cite{Li2024SASTCC, Dalaq2025SASTSLR}. Duas técnicas centrais sustentam a maior parte das análises relevantes para vulnerabilidades de injeção e de manipulação de dados sensíveis:


\begin{itemize}
\item \textbf{Análise sintática (baseada em padrões).} Ferramentas como o Bandit percorrem a AST do programa procurando por padrões estruturais conhecidos como inseguros ,por exemplo, a chamada de \texttt{subprocess} com \texttt{shell=True}, o uso de \texttt{eval()} sobre entrada não confiável, ou a construção de comandos SQL por concatenação de strings. Essa técnica é rápida e não exige modelagem de fluxo de dados, mas tende a gerar tanto falsos positivos (quando o padrão aparece em contexto seguro) quanto falsos negativos (quando a mesma falha é expressa de forma sintaticamente diferente da regra cadastrada).
\item \textbf{Análise de fluxo de dados e \textit{taint analysis}.} Ferramentas baseadas em consultas mais expressivas, como o Semgrep e o CodeQL, modelam o fluxo de valores "contaminados" (\textit{tainted}) desde uma fonte não confiável (\textit{source} , por exemplo, entrada do usuário, argumentos de linha de comando, leitura de arquivo) até um ponto sensível de uso (\textit{sink} , por exemplo, execução de comando, consulta SQL, desserialização), sinalizando um alerta quando esse fluxo ocorre sem passar por uma função de sanitização reconhecida (\textit{sanitizer}) \cite{Li2024SASTCC, Dalaq2025SASTSLR}. Essa técnica é mais precisa para vulnerabilidades de injeção que atravessam múltiplas linhas ou funções, mas seu custo computacional é maior e sua eficácia depende diretamente da completude do catálogo de \textit{sources}, \textit{sinks} e \textit{sanitizers} configurado para a linguagem-alvo.
\end{itemize}

\subsubsection{Vantagens e limitações}\label{subsubsec:sast-vantagens-limitacoes}

Por não exigir execução, o SAST pode ser aplicado desde as fases mais iniciais do ciclo de desenvolvimento ,inclusive sobre código incompleto ou não compilável  cobrir a totalidade do código analisável em uma única execução e apresentar comportamento determinístico e de baixo custo computacional, características que o tornam adequado à integração em pipelines de CI/CD \cite{Dalaq2025SASTSLR, Tamberg2025LLMHarness}.

Suas limitações, contudo, são estruturais e bem documentadas empiricamente, e não meramente incidentais. A primeira delas é a dependência direta da qualidade e da abrangência do conjunto de regras: uma dissecção aprofundada de ferramentas baseadas em consultas (CodeQL, Semgrep) mostrou que a maior parte dos falsos negativos observados decorre de consultas frágeis ou incompletas, e não de limitações do motor de análise em si \cite{Li2024SASTCC}. A segunda é a dificuldade em capturar vulnerabilidades que dependem de lógica de negócio, de configuração externa ou de fluxo interprocedural profundo, cenário em que a análise puramente sintática ou de fluxo local tende a falhar \cite{Zhu2024SASTAndroid}. A terceira é o custo de triagem imposto pelo volume de falsos positivos: um estudo longitudinal com seis ferramentas aplicadas a 116 repositórios C++ observou que a maioria dos alertas de algumas delas permanecia sem qualquer tratamento ao longo de anos, evidenciando baixa acionabilidade prática \cite{Aloraini2019SASTWarnings}.

Um padrão recorrente nas avaliações empíricas recentes é a combinação de precisão relativamente alta com \textit{recall} baixo , o SAST tende a "acertar quando alerta", mas a deixar passar despercebida a maioria das vulnerabilidades reais existentes no código analisado. Em uma comparação controlada de oito ferramentas sobre o benchmark sintético Juliet, mesmo a melhor ferramenta atingiu \textit{recall} médio de aproximadamente 0,58, restrito apenas às categorias de vulnerabilidade que ela declarava cobrir \cite{Esposito2024SASTComparison}. Esse resultado é relevante porque contraria a percepção difundida de que o principal problema do SAST seriam exclusivamente os falsos positivos: a cobertura limitada (falsos negativos) mostra-se igualmente crítica, e possivelmente mais grave em termos de risco residual \cite{Esposito2024SASTComparison, Li2023SASTJava}.

Essas limitações técnicas se traduzem em barreiras de adoção documentadas na indústria. Uma revisão de literatura cobrindo 89 estudos identificou os falsos positivos como a barreira mais citada à adoção de SAST, seguida da baixa qualidade das saídas, do tempo de configuração inicial e do esforço manual de correção \cite{Wadhams2024SASTBarriers}. Em linha semelhante, um estudo qualitativo com 20 profissionais da indústria constatou que a maioria enfrenta dificuldade em diagnosticar os resultados das ferramentas e demanda informações mais finas , como rastreamento do defeito e interpretação da causa , em vez de apenas métricas agregadas \cite{Li2025IndustrySAST}. Um relato de experiência sobre integração de SAST em fluxos de CI/CD, incluindo projetos Python/Django, também identificou notificações pouco claras como obstáculo prático relevante \cite{Chembakottu2025SASTUsability}. Adicionalmente, um estudo com 815 \textit{commits} que introduziram vulnerabilidades mostrou que a combinação de cinco ferramentas SAST sinalizava 78\% dos casos em nível de função ou arquivo, mas grande parte dos alertas individuais era irrelevante para a vulnerabilidade em questão, exigindo triagem humana significativa \cite{Charoenwet2024SecureReview}. Em conjunto, esses resultados indicam que a utilidade prática dos relatórios como clareza, explicação, priorização  constitui uma dimensão de avaliação tão relevante quanto as métricas clássicas de detecção, o que fundamenta a componente qualitativa deste trabalho .

Por fim, estudos que compararam abordagens baseadas em regras com abordagens baseadas em aprendizado de máquina observam capacidades complementares entre os dois paradigmas, mas alertam que combinações ingênuas a simples união dos resultados de ambos e podem elevar o \textit{recall} ao custo de degradar significativamente a precisão global \cite{Croft2021RuleLearningSAST}. Esse é um alerta metodológico diretamente relevante para o desenho da abordagem híbrida deste trabalho, que não deve assumir a priori que combinar SAST e IA produz ganhos automáticos.


\subsection{Ferramentas SAST avaliadas}\label{subsec:ferramentas-sast}
No ecossistema Python, o presente trabalho seleciona duas ferramentas SAST gratuitas e amplamente adotadas, representativas de estratégias técnicas distintas: uma análise sintática específica para a linguagem (Bandit) e uma análise baseada em correspondência de padrões multi-linguagem (Semgrep Community Edition). Uma revisão sistemática recente sobre avaliação de ferramentas SAST aponta que nenhuma ferramenta domina todas as categorias de vulnerabilidade, recomendando justamente o uso combinado de abordagens complementares \cite{Dalaq2025SASTSLR} , critério que orientou essa escolha.


\subsubsection{Bandit}\label{subsubsec:bandit}
O Bandit é um analisador estático desenvolvido especificamente para Python, que percorre a AST do código-fonte aplicando um conjunto de \textit{plugins} (regras) predefinidos, cada um associado a um identificador de teste e, na maioria dos casos, a uma categoria CWE \cite{Dalaq2025SASTSLR}. Por operar diretamente sobre a sintaxe da linguagem, o Bandit não exige compilação e é tipicamente rápido, mas sua cobertura está limitada às regras nativas do projeto, sem suporte nativo a análise de fluxo de dados interprocedural complexa.

As evidências quantitativas disponíveis sugerem que o Bandit, mesmo em configuração padrão ("de fábrica"), é comparativamente mais efetivo do que outras ferramentas gratuitas para falhas típicas do ecossistema Python. Em um benchmark com 108 vulnerabilidades reais extraídas de pacotes PyPI, o Bandit em configuração padrão obteve sensibilidade de aproximadamente 59,7\%, ainda que acompanhada de especificidade baixa; com regras customizadas construídas por um especialista, esse valor subiu para a faixa de 72–81\%, evidenciando forte dependência do conjunto de regras , o mesmo padrão observado para ferramentas baseadas em consultas na Seção ~\ref{subsubsec:sast-vantagens-limitacoes} \cite{Marquer2026SAGA}. Em um estudo sobre correção de código Python gerado por LLMs, o Bandit detectou entre 28\% e 46\% das amostras vulneráveis, taxa de duas a três vezes superior à do CodeQL no mesmo conjunto de dados, o que reforça sua adequação como detector de referência para o ecossistema Python \cite{Alrashedy2025FDSP}.

\subsubsection{Semgrep Community Edition}\label{subsubsec:semgrep-ce}
O Semgrep é uma ferramenta de análise estática multi-linguagem baseada em correspondência de padrões sintáticos (\textit{pattern matching}) sobre a AST, com suporte a regras customizadas escritas em uma sintaxe declarativa próxima ao próprio código-fonte da linguagem-alvo, o que reduz a barreira de criação de novas regras em comparação a linguagens de consulta mais especializadas \cite{Dalaq2025SASTSLR}. A edição \textit{Community} (gratuita, de código aberto) inclui um conjunto de regras públicas mantidas pela comunidade e pela empresa mantenedora, sem os recursos de análise interprocedural mais avançada disponíveis nas edições comerciais.

O desempenho do Semgrep evidenciado na literatura reforça, mais uma vez, a forte dependência do conjunto de regras utilizado. No mesmo benchmark de 108 vulnerabilidades reais de pacotes PyPI citado acima, o Semgrep em configuração padrão obteve sensibilidade de apenas cerca de 28,6\% \cite{Marquer2026SAGA}. Em Java, um estudo com 170 \textit{commits} associados a CVEs mostrou que a taxa de detecção do Semgrep quase triplicou de 15,3\% para 44,7\% , após engenharia manual de regras específicas para o conjunto avaliado \cite{Bennett2024Semgrep}. Esse resultado, obtido fora do ecossistema Python, é relevante para o TCC como alerta metodológico: o desempenho do Semgrep Community Edition em configuração padrão, que é a configuração adotada neste experimento por representar o que está prontamente disponível a um desenvolvedor sem investimento adicional em engenharia de regras, deve ser interpretado como um limite inferior do potencial da ferramenta.

Outras ferramentas do ecossistema , notadamente o CodeQL e o analisador de fluxo de dados Pysa  foram consideradas, mas não compõem o braço SAST deste experimento. O CodeQL, em uma análise longitudinal sobre CVEs reais (incluindo projetos Python), apresentou taxa de detecção agregada de apenas 4,4\%, além de regressões de detecção entre versões sucessivas da ferramenta \cite{Noirot2025CodeQLLongitudinal}; o Pysa, no benchmark de pacotes PyPI já citado, apresentou sensibilidade inferior a 1\% em configuração padrão \cite{Marquer2026SAGA}. Bandit e Semgrep Community Edition foram, portanto, priorizados por serem gratuitos, amplamente adotados na prática e representativos de duas estratégias técnicas distintas e complementares.


\section{Modelos de Linguagem de Grande Escala (LLMs)}\label{sec:llms}
\subsection{Conceitos fundamentais}\label{subsec:llms-conceitos}
\subsubsection{Arquitetura Transformer}\label{subsubsec:transformer}
Os LLMs relevantes para este trabalho - os modelos por trás de assistentes de programação como GitHub Copilot, e das ferramentas empresariais avaliadas no experimento, Cursor e OpenAI Codex — são construídos sobre a arquitetura \textbf{Transformer}, proposta por \citeonline{Vaswani2017Transformer}. Sem aprofundar a formulação matemática, que foge ao escopo deste trabalho, cabe destacar a característica central da arquitetura: em vez de processar uma sequência de texto (ou de código) de forma recorrente, token a token, o Transformer utiliza um mecanismo de \textbf{atenção} (\textit{self-attention}) que permite a cada elemento da sequência ponderar diretamente sua relação com todos os demais elementos, independentemente da distância entre eles. Essa propriedade viabilizou o treinamento eficiente e paralelizável de modelos com bilhões de parâmetros e é a base arquitetural comum a praticamente todos os LLMs de propósito geral e especializados em código utilizados atualmente, incluindo os modelos que sustentam as ferramentas avaliadas neste TCC.

\subsubsection{Treinamento, tokenização e janela de contexto}\label{subsubsec:tokenizacao-contexto}
O treinamento de um LLM ocorre tipicamente em duas etapas principais: um \textbf{pré-treinamento} auto-supervisionado sobre grandes volumes de texto (e, no caso de modelos de código, também código-fonte), em que o modelo aprende a prever o próximo token de uma sequência, seguido de etapas de ajuste fino (\textit{fine-tuning}) e alinhamento , por exemplo, via aprendizado por reforço com retroalimentação humana , voltadas a tornar o modelo mais útil e seguro em interações reais. Antes de ser processado pelo modelo, o texto de entrada é convertido em uma sequência de \textbf{tokens} que são unidades que podem corresponder a palavras inteiras, subpalavras ou símbolos de código, conforme o vocabulário do tokenizador utilizado e é sobre essa sequência de tokens que o mecanismo de atenção.

A \textbf{janela de contexto} (\textit{context window}) é o número máximo de tokens que o modelo pode processar simultaneamente em uma única interação, somando a instrução do usuário, o histórico da conversa e todo o conteúdo de código fornecido como contexto. Essa limitação é diretamente relevante para a análise de vulnerabilidades em nível de repositório: quanto maior a janela de contexto, maior a quantidade de código-fonte (múltiplos arquivos, dependências, histórico de chamadas de função) que pode ser fornecida ao modelo em uma única análise. A literatura recente, porém, mostra que uma janela de contexto ampla não elimina o problema de utilização eficiente desse contexto , visto que modelos degradam sua capacidade de localizar informações relevantes conforme o comprimento do texto de entrada aumenta, fenômeno descrito na literatura como \textit{lost-in-the-end} no contexto específico de detecção de vulnerabilidades \cite{Sovrano2025LostInEnd}.


\subsection{LLMs aplicados a código}\label{subsec:llms-codigo}
\subsubsection{Pré-treinamento em código-fonte}\label{subsubsec:pre-treinamento-codigo}
A aplicação de LLMs a tarefas de programação foi viabilizada pelo pré-treinamento (ou pré-treinamento contínuo) sobre grandes volumes de código-fonte público, tipicamente extraído de repositórios abertos. O marco histórico dessa linha é o \textbf{Codex}, um GPT ajustado sobre código do GitHub e avaliado quanto à correção funcional de programas gerados a partir de descrições em linguagem natural; o Codex é o modelo que originou o GitHub Copilot e, por extensão, deu nome à atual ferramenta "OpenAI Codex" avaliada neste trabalho \cite{Chen2021Codex}. Na mesma linha histórica, modelos abertos como o \textbf{Code Llama}, derivado do Llama 2 da Meta AI via pré-treinamento contínuo em código, com variantes especializadas em preenchimento de lacunas (\textit{infilling}) e em janelas de contexto estendidas, ilustram a consolidação do paradigma de modelos de propósito específico para código \cite{Roziere2023CodeLlama}. Esses modelos são citados neste trabalho como contexto histórico e técnico, e não como objeto de avaliação empírica direta — o experimento deste TCC avalia as ferramentas empresariais construídas sobre modelos dessa linhagem, não os modelos isoladamente.


\subsubsection{Capacidades e limitações}\label{subsubsec:llms-capacidades-limitacoes}
A literatura documenta capacidades relevantes de LLMs em tarefas relacionadas a código , como geração, explicação, transformação e, de interesse direto para este trabalho, identificação de padrões potencialmente inseguros , acompanhadas de limitações sistemáticas que precisam ser consideradas na interpretação de qualquer avaliação empírica. A principal delas é a \textbf{alucinação}: a geração de afirmações plausíveis, mas factualmente incorretas ou não fundamentadas no código realmente fornecido, fenômeno que no contexto de detecção de vulnerabilidades se manifesta como falsos positivos com justificativas aparentemente coerentes, mas tecnicamente equivocadas. A segunda é a \textbf{dependência de prompt}: pequenas variações na formulação da instrução produzem diferenças substanciais de desempenho. Estudos citados no Capítulo ~\ref{cap:estadoarte} relatam ganhos de até 0,18 em F1-score entre estratégias de \textit{prompt} \cite{Khare2025LLMEffectiveness} e diferenças superiores a 0,3 em F1-score entre formulações de instrução \cite{Iranmanesh2025ZeroFalse} , o que compromete a reprodutibilidade das avaliações se o protocolo experimental não fixar e documentar cuidadosamente os \textit{prompts} utilizados. A terceira é a \textbf{previsibilidade limitada} do comportamento do modelo entre execuções: mesmo sob a mesma entrada, LLMs comerciais apresentam variação de resultado entre chamadas idênticas, um comportamento não-determinístico documentado já nas primeiras comparações entre LLMs e SAST \cite{Ozturk2023AIvsSAST}, discutido com mais profundidade no Capítulo ~\ref{cap:estadoarte}.

\subsection{Retrieval-Augmented Generation (RAG)}\label{subsec:rag}
Um LLM de propósito geral, por si só, tem acesso apenas ao conhecimento codificado em seus parâmetros durante o treinamento, o que inclui uma amostra de código público, mas não o conteúdo específico do repositório privado que se deseja analisar. O paradigma de \textbf{Retrieval-Augmented Generation (RAG)}, formalizado por \citeonline{Lewis2020RAG}, propõe combinar o modelo generativo com um componente de recuperação (\textit{retriever}) que busca, em uma base externa, os trechos mais relevantes para a tarefa em questão, fornecendo-os ao modelo como contexto adicional antes da geração da resposta. O objetivo é reduzir a alucinação e permitir que o modelo responda com base em informação concreta e verificável, em vez de depender exclusivamente do que foi memorizado durante o treinamento.

Esse princípio é diretamente relevante para entender como as ferramentas de codificação avaliadas neste trabalho operam sobre um repositório: em vez de receber todo o código-fonte de uma vez (o que seria inviável para repositórios grandes, dada a limitação de janela de contexto), essas ferramentas implementam mecanismos de busca por indexação semântica de arquivos, busca textual, navegação guiada pela estrutura do projeto para localizar e fornecer ao modelo apenas os trechos de código considerados relevantes para a tarefa corrente. É importante, contudo, não superestimar essa analogia: ferramentas comerciais de codificação não implementam necessariamente RAG na forma canônica descrita por \citeonline{Lewis2020RAG} , mas sim variações e heurísticas proprietárias inspiradas no mesmo princípio geral de fundamentar a geração em contexto recuperado. Essa distinção é relevante porque a forma como o contexto é entregue ao modelo afeta diretamente seu desempenho em tarefas de análise de segurança: avaliações citadas no Capítulo ~\ref{cap:estadoarte} mostram que abordagens que fornecem contexto estruturado de fatias de código guiadas por grafos de propriedade \cite{Lekssays2025LLMxCPG}, abstrações de funções chamadas \cite{Yang2025ContextVD} ou especificações extraídas de histórico de correções \cite{Zhu2025SpecGuidedVD} , melhoram consistentemente os resultados de detecção em relação a análises de trechos de código isolados, o que reforça a relevância de avaliar ferramentas com acesso real ao repositório, e não apenas modelos respondendo sobre um trecho isolado de código.


\section{Agentes de Codificação Assistidos por IA}\label{sec:agentes-codificacao}
\subsection{De autocomplete a agentes: evolução do paradigma}\label{subsec:evolucao-agentes}
O uso de LLMs no desenvolvimento de software evoluiu em estágios sucessivos de autonomia. Na primeira geração, ferramentas como o GitHub Copilot atuavam como \textbf{completadores de código} (\textit{autocomplete}): a partir do contexto imediato (o texto já digitado no arquivo), o modelo sugeria a continuação mais provável de uma linha ou um pequeno bloco , e cada sugestão exigia aceitação explícita do desenvolvedor antes de ser incorporada ao código. Nessa fase, o LLM não tinha iniciativa: reagia a um contexto local sem planejar múltiplos passos nem interagir com o ambiente além da própria sugestão de texto.

Uma segunda geração introduziu assistentes conversacionais integrados ao ambiente de desenvolvimento, capazes de gerar blocos de código mais extensos, ou de propor alterações em múltiplos trechos de um arquivo, a partir de instruções em linguagem natural .Ainda, contudo, sob iniciativa do desenvolvedor a cada interação. O terceiro estágio, no qual se inserem as ferramentas avaliadas neste trabalho, é o de \textbf{agentes de codificação}: sistemas em que o modelo decompõe uma tarefa de alto nível em subtarefas, decide autonomamente quais ferramentas invocar (busca no repositório, leitura e escrita de arquivos, execução de comandos de terminal, execução de testes ou de analisadores), avalia os resultados obtidos e itera até considerar a tarefa concluída , com supervisão humana em pontos de controle definidos, mas sem exigir aprovação a cada micro-ação. A consolidação desse paradigma em engenharia de software foi impulsionada por benchmarks de resolução autônoma de \textit{issues} reais de repositório, como o SWE-bench \cite{Jimenez2024SWEbench}, e pela demonstração de que a qualidade da interface entre o agente e o ambiente computacional impacta fortemente seu desempenho \cite{Yang2024SWEagent}.


\subsection{O que caracteriza um ``agente''}\label{subsec:caracterizacao-agente}
O ciclo de operação de um agente de LLM é formalizado, na literatura, pelo paradigma \textbf{ReAct} (\textit{Reasoning and Acting}): o modelo intercala passos explícitos de raciocínio (\textit{thought})  em que articula, em linguagem natural, o que precisa ser feito e por quê , com ações concretas (\textit{action}) e invocações de ferramentas externas ao próprio modelo , além de observações (\textit{observation}) do resultado dessas ações no ambiente, repetindo esse ciclo até concluir a tarefa ou atingir um critério de parada \cite{Yao2023ReAct}. Três elementos, tomados em conjunto, distinguem um agente de um simples completador de texto ou de um chatbot de pergunta-e-resposta:

1. \textbf{Uso de ferramentas (\textit{tool-calling}).} O modelo pode invocar funções externas, como leitura e escrita de arquivos, execução de comandos de shell, chamadas a analisadores estáticos ou a interpretadores  e incorporar o resultado dessas chamadas ao seu raciocínio subsequente, em vez de responder apenas com base em seu conhecimento paramétrico.

2. \textbf{Iteração autônoma.} O modelo decide, por conta própria, quantos ciclos de raciocínio-ação são necessários para concluir a tarefa, sem que cada passo intermediário exija uma nova instrução explícita do usuário.

3. \textbf{Acesso e leitura do estado do ambiente.} O agente pode explorar a estrutura de um repositório, ler múltiplos arquivos relacionados e navegar entre eles conforme a necessidade identificada durante o próprio raciocínio, em vez de depender exclusivamente do trecho de código fornecido a priori pelo usuário.

A formalização de \textbf{interfaces agente-computador} (\textit{Agent-Computer Interfaces}) , o conjunto de ferramentas e formatos de retorno disponibilizados ao agente demonstrou que a qualidade dessa interface afeta diretamente o desempenho observável: com uma interface especializada, a taxa de resolução de \textit{issues} reais no benchmark SWE-bench passou de 1,31\% (quando o modelo apenas recuperava contexto de forma simples) para 12,47\% \cite{Yang2024SWEagent, Jimenez2024SWEbench}. Ainda assim, as taxas absolutas de resolução permanecem baixas em tarefas gerais de engenharia de software, o que calibra expectativas realistas sobre a autonomia desses sistemas em tarefas complexas de repositório \cite{Jimenez2024SWEbench} , calibração especialmente relevante para a interpretação dos resultados deste TCC, cujo domínio (detecção de vulnerabilidades) é tipicamente mais difícil de verificar automaticamente do que a resolução de um \textit{bug} funcional com um conjunto de testes associado.

\subsection{Modelos, agentes e camada operacional}\label{subsec:camada-operacional-agente}
Para interpretar corretamente o objeto de avaliação deste trabalho, é necessário distinguir três níveis. O \textbf{LLM isolado} recebe uma entrada delimitada e produz texto, uma classificação ou chamadas de ferramenta; por si só, não explora um repositório nem executa ações no ambiente. O \textbf{agente de codificação} organiza ciclos de raciocínio, ação e observação para cumprir uma tarefa, conforme descrito na Seção ~\ref{subsec:caracterizacao-agente}. Por fim, a \textbf{camada operacional do agente} conecta o modelo ao ambiente: seleciona e recupera contexto, mantém o estado da tarefa, disponibiliza ferramentas, executa comandos, aplica permissões e define critérios de parada. A literatura recente também denomina essa camada de \textit{agent harness}, embora o termo ainda não seja plenamente padronizado \cite{Yang2024SWEagent, Macedo2026AgentHarness}.

Cursor e OpenAI Codex não são, portanto, modelos de IA isolados nem devem ser tratados como sinônimo de "IA pura". São ferramentas empresariais de desenvolvimento assistido por IA que implementam agentes de codificação apoiados por uma camada operacional própria. Seus resultados refletem a configuração completa ferramenta-agente onde contém modelo(s), recuperação de contexto, ferramentas disponíveis, permissões e iteração  e não somente uma capacidade intrínseca de um modelo de linguagem. Neste trabalho, a expressão \textbf{LLM isolado} é reservada aos estudos em que o modelo recebe código em uma entrada fixa, sem exploração instrumental autônoma do repositório; Cursor e Codex constituem o braço de \textbf{ferramentas-agente empresariais}.

\subsection{Ferramentas avaliadas}\label{subsec:ferramentas-agentes}
Este trabalho avalia duas ferramentas empresariais de desenvolvimento assistido por IA que implementam o paradigma de agente descrito nas Seções ~\ref{subsec:evolucao-agentes} e ~\ref{subsec:caracterizacao-agente}, operando com acesso direto ao repositório de código sob análise.

\subsubsection{Cursor}\label{subsubsec:cursor}
O Cursor é um editor de código construído sobre a base do Visual Studio Code, que incorpora um modo de operação agêntico (\textit{Agent Mode}) no qual o modelo compreende a estrutura do repositório, propõe e executa alterações em múltiplos arquivos, executa comandos de terminal e itera sobre o resultado dessas ações até considerar a tarefa concluída, mediante instruções em linguagem natural \cite{Cursor2025Docs}. No contexto deste trabalho, o Cursor é operado em modo agente para realizar a análise de segurança do código Python fornecido, com acesso de leitura ao repositório e, conforme o protocolo experimental definido na metodologia, possivelmente também a ferramentas de execução para validação de hipóteses sobre o comportamento do código.

\subsubsection{OpenAI Codex}\label{subsubsec:openai-codex}
O OpenAI Codex, na configuração avaliada neste trabalho, refere-se ao agente de codificação de linha de comando (\textit{Codex CLI}) disponibilizado pela OpenAI, que executa localmente com acesso configurável a arquivos, comandos de terminal e políticas de permissão (\textit{sandbox}) \cite{OpenAI2025CodexDocs}. É importante notar a ambiguidade histórica do nome: "Codex" designou originalmente o modelo de linguagem ajustado para código descrito por \citeonline{Chen2021Codex}, hoje descontinuado como produto isolado; a ferrasmenta atual reutiliza o nome, mas opera como um agente autônomo construído sobre modelos mais recentes, e não como o completador de código original. Essa distinção é registrada aqui para evitar confusão entre o objeto de estudo empírico deste TCC (a ferramenta-agente atual) e a referência histórica citada na Seção ~\ref{subsubsec:pre-treinamento-codigo}.


\subsection{Uso de IA para segurança de código}\label{subsec:ia-seguranca-codigo}
\subsubsection{Detecção de vulnerabilidades via LLM}\label{subsubsec:deteccao-llm}
Para além da geração de código, LLMs vêm sendo empregados diretamente como detectores de vulnerabilidades, seja por meio de \textit{prompts} que solicitam a análise de um trecho de código, seja por meio de agentes com acesso ao repositório. Surveys recentes documentam a rápida expansão dessa linha de pesquisa, catalogando de 58 a 208 estudos primários sobre o tema \cite{Sheng2025LLMSurvey, Zhou2025LLMReview, Germano2025LLMSLR}. O padrão geral observado que será detalhado com maior profundidade no Capítulo ~\ref{cap:estadoarte} , é de potencial de cobertura mais amplo que o das ferramentas SAST tradicionais, acompanhado de taxas de falso positivo mais altas, sensibilidade a variações de \textit{prompt} e de benchmark, e desempenho substancialmente reduzido quando avaliado sob condições metodologicamente mais rigorosas (dados sem vazamento entre treino e teste, comparações em pares entre versão vulnerável e corrigida) \cite{Ding2025CodeLMVuln, Khare2025LLMEffectiveness}.

Estudos de campo indicam, adicionalmente, uma distância relevante entre o desempenho reportado em benchmark e a utilidade percebida por desenvolvedores reais: em um estudo com 17 desenvolvedores utilizando uma ferramenta de detecção e reparo baseada em IA integrada ao ambiente de desenvolvimento, os alertas gerados receberam utilidade média de 2,5 em uma escala de 5 pontos, e 75\% das correções sugeridas não eram diretamente aplicáveis \cite{Steenhoek2025IDEUserStudy}. Esse resultado reforça, no contexto específico de ferramentas de IA, a mesma conclusão já observada para SAST na Seção ~\ref{subsubsec:sast-vantagens-limitacoes}: métricas de detecção favoráveis não implicam, automaticamente, relatórios úteis na prática de desenvolvimento — o que fundamenta a avaliação qualitativa conduzida neste trabalho (Seção ~\ref{subsec:metricas-qualitativas}).

\subsubsection{Abordagens híbridas SAST + IA}\label{subsubsec:hibrido-sast-ia}
Diante do padrão de \textit{trade-off} entre SAST (preciso, porém de cobertura limitada) e LLMs (cobertura potencialmente maior, porém com mais falsos positivos e menor determinismo), uma linha de pesquisa crescente propõe combinar as duas abordagens. O padrão arquitetural mais próximo do desenho experimental deste TCC posiciona o LLM a jusante do SAST: a ferramenta estática gera os alertas iniciais, e o LLM os revisa, filtra, explica ou valida com base no contexto do código, atuando como uma camada de triagem inteligente. Resultados reportados nessa linha incluem reduções substanciais de falsos positivos mantendo a maior parte dos verdadeiros positivos \cite{Ameen2026QASecClaw, Iranmanesh2025ZeroFalse, Agrawal2025SASTGenius}, ainda que nem sempre de forma uniforme entre categorias de vulnerabilidade — em alguns cenários, a filtragem por LLM chega a piorar o desempenho em categorias específicas, e o SAST puro permanece preferível quando o \textit{recall} é prioridade absoluta \cite{Ciavatta2026LLMSAST}.

Outros padrões arquiteturais posicionam o LLM a montante do motor estático , por exemplo, inferindo especificações de \textit{taint} (pares \textit{source}/\textit{sink}) para alimentar ferramentas como o CodeQL \cite{Li2025IRIS} . Ou utilizam o SAST como oráculo de realimentação em laços iterativos de geração e correção de código, abordagem já demonstrada especificamente em Python \cite{Alrashedy2025FDSP, Lee2026CodeEnhancer}. Esses três padrões de filtragem a jusante, ampliação a montante e realimentação em laço  são discutidos com maior detalhamento empírico, incluindo as ressalvas metodológicas sobre a maturidade ainda incipiente dessa literatura (concentração em preprints, escassa validação em Python), no Capítulo ~\ref{cap:estadoarte}.


\section{Métricas de Avaliação de Ferramentas de Detecção}\label{sec:metricas-avaliacao}
\subsection{Métricas quantitativas}\label{subsec:metricas-quantitativas}
A avaliação quantitativa de detectores de vulnerabilidades , sejam eles baseados em regras (SAST) ou em LLMs, utiliza majoritariamente métricas derivadas da matriz de confusão: \textbf{precisão} (proporção de alertas emitidos que correspondem a vulnerabilidades reais), \textbf{recall} ou sensibilidade (proporção de vulnerabilidades reais efetivamente detectadas), \textbf{F1-score} (média harmônica entre precisão e recall) e a \textbf{taxa de falsos positivos} \cite{Kalouptsoglou2023VulnPrediction, Esposito2024SASTComparison}. Um mapeamento sistemático com 180 estudos identificou o F1-score como a métrica mais popular, mas recomendou o uso complementar de variantes que penalizam mais fortemente os falsos negativos, como o \textbf{F2-score}, sob o argumento de que, em segurança, deixar de detectar uma vulnerabilidade real tende a ser mais grave do que gerar um alerta espúrio que será descartado na triagem \cite{Kalouptsoglou2023VulnPrediction}. Na mesma direção, benchmarks recentes adotam o \textbf{F3-score} como métrica primária, penalizando ainda mais os falsos negativos \cite{Pellew2026RealVuln}, e a revisão sistemática de \citeonline{Dalaq2025SASTSLR} recomenda o \textbf{índice de Youden} como medida balanceada entre sensibilidade e especificidade, menos sensível ao desbalanceamento de classes do que o F1 isolado.

A literatura recente alerta, contudo, que métricas agregadas podem ser enganosas quando calculadas sobre benchmarks metodologicamente frágeis. Foi demonstrado que rótulos incorretos e vazamento de dados entre conjuntos de treino e teste inflam artificialmente os resultados reportados: um modelo com F1 de 68,3\% em um benchmark tradicional caiu para 3,1\% quando reavaliado em um conjunto de dados corrigido e deduplicado \cite{Ding2025CodeLMVuln}. Como resposta a esse problema, foram propostas métricas mais rigorosas, como o \textbf{VD-S} (taxa de falsos negativos sob uma taxa de falsos positivos limitada) e a \textbf{avaliação em pares} (comparar o julgamento do detector sobre a função vulnerável e sobre sua versão corrigida), que penaliza detectores que simplesmente classificam tudo como vulnerável, prática que infla artificialmente o \textit{recall} aparente \cite{Ding2025CodeLMVuln, Yildiz2025LLMAgents}.


\subsection{Métricas qualitativas}\label{subsec:metricas-qualitativas}
Métricas puramente quantitativas não capturam integralmente a utilidade de um relatório de vulnerabilidade para um desenvolvedor em contexto real de trabalho, como evidenciado repetidamente ao longo deste capítulo (Seções ~\ref{subsubsec:sast-vantagens-limitacoes} e ~\ref{subsubsec:deteccao-llm}). Cresce, por isso, o reconhecimento de que a utilidade prática dos relatórios deve ser avaliada explicitamente, por meio de dimensões qualitativas complementares:

\begin{itemize}
\item \textbf{Clareza da explicação} — se o relatório comunica, em linguagem compreensível, por que o trecho de código é considerado vulnerável, e não apenas qual regra ou padrão foi acionado.
\item \textbf{Qualidade da recomendação de correção} — se a sugestão fornecida (quando existente) é tecnicamente correta, específica ao contexto do código e diretamente aplicável, ou genérica a ponto de exigir retrabalho substancial do desenvolvedor.
\item \textbf{Indicação de severidade e priorização} — se o relatório ajuda a hierarquizar múltiplos achados por criticidade, facilitando a triagem em cenários com grande volume de alertas, precisamente o gargalo identificado como principal barreira de adoção do SAST na Seção ~\ref{subsubsec:sast-vantagens-limitacoes}.
\end{itemize}
Rubricas qualitativas formais foram propostas para operacionalizar essa avaliação, notadamente a rubrica de \textbf{corretude, clareza e acionabilidade} (\textit{accuracy, clarity, actionability}) \cite{Mao2025ExplainableVD}, que constatou que uma alta taxa de detecção não implica, por si só, explicações acionáveis. Estudos com profissionais da indústria corroboram essa dissociação: alertas tecnicamente corretos podem ser percebidos como pouco úteis no fluxo real de trabalho \cite{Steenhoek2025IDEUserStudy, Li2025IndustrySAST}. Este trabalho combina, portanto, métricas quantitativas (precisão, \textit{recall}, F1, taxa de falsos positivos) com avalfiação qualitativa dos relatórios produzidos pelos três braços do experimento (clareza, qualidade da recomendação e indicação de severidade), seguindo esse consenso metodológico emergente na literatura.

\subsection{Benchmarks de vulnerabilidades}\label{subsec:benchmarks-vulnerabilidades}
A validade de qualquer comparação quantitativa entre ferramentas depende diretamente da qualidade do benchmark utilizado , isto é, do conjunto de código rotulado como vulnerável ou seguro sobre o qual as métricas da Seção ~\ref{subsubsec:bandit} são calculadas. A literatura recente demonstra que benchmarks amplamente utilizados sofrem de problemas metodológicos sistemáticos: rótulos incorretos, vazamento de dados entre partições de treino e teste, e distribuições de vulnerabilidade pouco realistas em relação a código de produção, todos fatores que inflam artificialmente os resultados reportados \cite{Ding2025CodeLMVuln, Chen2023DiverseVul, Chakraborty2022DLLimits}. Avaliações conduzidas em nível de função isolada tendem, adicionalmente, a superestimar a capacidade real dos detectores em comparação a avaliações em nível de repositório completo, onde o contexto necessário para julgar corretamente um trecho de código está frequentemente disperso em múltiplos arquivos \cite{Wen2026VulEval}.

Diante desse cenário, benchmarks mais recentes e cuidadosamente curados assumem papel central na avaliação empírica deste trabalho, com destaque para o \textbf{RealVuln}, benchmark focado especificamente no ecossistema Python, com dados rotulados a partir de vulnerabilidades reais em repositórios de código aberto \cite{Pellew2026RealVuln}. Sua metodologia de construção, seus resultados de comparação entre ferramentas de regras, LLMs de propósito geral e \textit{scanners} especializados, e sua relação direta com o protocolo experimental adotado neste TCC são detalhados no Capítulo ~\ref{cap:estadoarte}, no contexto mais amplo do estado da arte em comparações entre SAST e abordagens baseadas em IA.
